%%%%%%%%%%  scan idx 173  |  folio 162  %%%%%%%%%%

2.3. \underline{Morfismos de biextensiones}. Sea $E$ una biextensión
de $(P,Q)$ por $G$, y sea asimismo $E'$ una biextensión de $(P',Q')$
por $G'$, donde $P',Q',G'$ son también grupos conmutativos de
$\underline T$. Se llama \underline{morfismo de biextensiones} de $E$
a $E'$ a un sistema $(u,v,w,f)$, donde
\begin{equation}\tag{2.3.1}
 u:P\longrightarrow P',\qquad
 v:Q\longrightarrow Q',\qquad
 w:G\longrightarrow G'
\end{equation}
son morfismos de grupos, y
\begin{equation}\tag{2.3.2}
 f:E\longrightarrow E'
\end{equation}
es un morfismo de los haces de conjuntos subyacentes a $E$ y $E'$, tal
que $f$ es simultáneamente un homomorfismo de extensiones de grupos
sobre $P$ y $P'$, respectivamente, asociado a los homomorfismos de
grupos relativos
\[
 u\times v:Q_P\longrightarrow Q'_{P'},
 \qquad
 u\times w:G_P\longrightarrow G'_{P'},
\]
y un homomorfismo de extensiones de grupos sobre $Q$ y $Q'$, asociado
a los homomorfismos de grupos relativos
\[
 u\times v:P_Q\longrightarrow P'_{Q'},
 \qquad
 v\times w:G_Q\longrightarrow G'_{Q'}.
\]
Obsérvese además que $u,v,w$ quedan determinados por $f$ solamente. En
efecto, pasando al cociente se recupera $u\times v$ a partir de $f$, y
después se recuperan $u$ y $v$ por restricción a $P\times1_Q$ y
$1_P\times Q$, respectivamente; restringiendo $f$ a $G.e_E$ se
recupera $w$. No hay, pues, inconveniente en identificar el morfismo
$(u,v,w,f)$ con $f$.

Con la precaución habitual de adjuntar a los morfismos sus fuentes y
sus metas (la genuflexión canónica), se ve que \underline{las
biextensiones del topos $\underline T$ son los objetos de una
categoría}, cuyos objetos y morfismos son los definidos en 2.1 y 2.3 y
cuya composición

%%%%%%%%%%  scan idx 174  |  folio 163  %%%%%%%%%%

de morfismos se define de manera evidente. Además,
$(P,Q,G,E)\mapsto(P,Q,G)$ da un funtor canónico
\begin{equation}\tag{2.3.3}
 \operatorname{BIEXT}(\underline T)
 \longrightarrow
 \operatorname{GRAB}(\underline T)\times
 \operatorname{GRAB}(\underline T)\times
 \operatorname{GRAB}(\underline T),
\end{equation}
donde $\operatorname{BIEXT}(\underline T)$ designa la categoría de
biextensiones sobre $\underline T$ y
$\operatorname{GRAB}(\underline T)$ la categoría de grupos abelianos
sobre $\underline T$. Conviene descomponer este funtor canónico en los
dos funtores $E\mapsto(P,Q)$,
\begin{equation}\tag{2.3.4}
 \operatorname{BIEXT}(\underline T)
 \longrightarrow
 \operatorname{GRAB}(\underline T)\times
 \operatorname{GRAB}(\underline T),
\end{equation}
y $E\mapsto G$,
\begin{equation}\tag{2.3.5}
 \operatorname{BIEXT}(\underline T)
 \longrightarrow
 \operatorname{GRAB}(\underline T).
\end{equation}

